\chapter{The atom is a point carrying values}

\section{One part, several magnitudes}

A part in chemistry is an atom, and an atom is not one number. It carries a count of protons \(Z\),
which names the element and orders the whole set; a count of valence electrons, the outer
shell's occupancy; a radius, in two readings, the covalent one at which it bonds and the van der
Waals one at which it merely touches; and an electronegativity, the pull it puts on a shared pair.
Written together these are a vector, and every reading in this book is a magnitude taken from it or a
distance taken between two of them. The engine reads any domain as points carrying values, and the
atom is that description with the values fixed by the element.

Two atoms in a molecule are two such points at two positions, and the bond between them is the
vector from one to the other. Its magnitude is the bond length. That is the whole of the
representation: positions carrying element vectors, and the differences between positions carrying
lengths. Nothing about it is particular to chemistry until the values are filled in. That is why the
part of the engine that knows an atom exists is the only part that knows any domain exists.

\section{The order is a real reading, not a numbering}

The elements are ordered by \(Z\), and that order is not a label chosen for convenience. Read a
property along it, the atomic radius or the first ionization energy, and the sequence does not
wander. It rises and falls with a period, and the period is the shell filling: two, then eight, then
eight, then eighteen, then eighteen, then thirty-two. The periodic law is the statement that a
property of the atom is a periodic function of \(Z\), and in the engine's terms it is the first
departure chemistry offers. A property read along \(Z\) agrees with itself at the shell boundaries and
disagrees with a shuffle of the same values.

The shell counts are themselves a departure from a flat expectation. A shell \(n\) holds \(2n^2\)
electrons because that is how many distinct states the quantum numbers allow it. The capacities
are \(2, 8, 18, 32\) and the period lengths follow from the order the sub-shells fill. This is not a
count anybody chose. It is fixed before any element is weighed, the property the later
chapters lean on: an answer known before the measurement and not supplied by it.

The element order, the electron set and the shell counts are the atomic structure, and this book
reads them as given. It does not establish them and does not transcribe them. They are the ledger
chemistry consumes, authored once where the atom is described, and everything this book adds is the
bonding that sits on top of that ledger.

\section{The set is not bounded}

There is a temptation to write chemistry as a fixed table of some number of elements. The engine's
record refuses it. Every bound put on a quantity in this work has had to be taken back off, and a
bounded element set is the same defect waiting in a new place. The count of known elements is a fact
about what has been synthesized, not about what the periodic law describes, and the law describes a
sequence with no last term until a nucleus cannot hold together. A reading that fixes the set at a
round number reads that number and not the chemistry.

The same holds one level down. An element is not one atom either. It is a set of isotopes differing
in neutron count, weighted by abundance, and the weighting is a measured distribution and not a
constant. Where a reading needs a mass it needs the distribution behind it, and a single tabulated
value is a summary of that distribution with its spread deleted. The building block is a set, at
every level at which it is read, and the engine's discipline is to draw the background from the set
and not to cap it.
